% Copyright 2026 The terCAD team. All rights reserved.
% Use of this content is governed by a CC BY-NC-ND 4.0 license that can be found in the LICENSE file.

\subsection{Configuring Environment} \label{first-step}
\markboth{Prototyping}{Configuring Environment}

Let's assume that we are starting from scratch with an empty folder and a new operating system.
The good news is that NativePHP doesn't require any exotic tools. The bad news is that it needs
all of them at once: a PHP toolchain to run Laravel, a Node.js toolchain to build front-end
assets, and native toolchains to compile Android and iOS shells. The following are the
requirements for NativePHP v4: (\href{https://nativephp.com/docs/mobile/4/getting-started/environment-setup}{https://nativephp.com/docs/mobile/4/getting-started/environment-setup}):

\begin{itemize}
  \item PHP 8.3+ and Laravel 10+, Composer, and Node.js with npm (for Vite).

  \item Android -- Android Studio 2024.2.1+, Android SDK with API 29+ (Build-Tools and
  Platform-Tools), and JDK (the documentation uses JDK 17 in examples). On Windows, 7-Zip is also
  required, and NativePHP doesn't run under WSL.

  \item iOS -- an Apple silicon Mac (M1+) with macOS 15.6+, Xcode 26+ with Command Line Tools,
  and CocoaPods. An Apple Developer account isn't needed for the simulator, but it's required
  for real devices and the App Store.
\end{itemize}

\noindent On Linux (Ubuntu-based distributions), PHP and its extensions are installed by the
package manager:

\begin{lstlisting}[language=terminal]
$ sudo apt install php8.4 php8.4-cli php8.4-mbstring php8.4-xml php8.4-zip php8.4-curl
$ sudo apt install zip unzip
# switch the default interpreter, if several versions are installed
$ sudo update-alternatives --set php /usr/bin/php8.4
\end{lstlisting}

\noindent Android Studio is downloaded from
\href{https://developer.android.com/studio}{https://developer.android.com/studio}. After the
initial wizard, the SDK Manager (Tools $\rightarrow$ SDK Manager) installs an SDK Platform and
SDK Tools. Then the terminal has to know where the SDK and JDK are located. It's the most
frequent reason for the first failed build:

\begin{lstlisting}[language=terminal]
$ echo 'export ANDROID_HOME=$HOME/Android/Sdk' >> ~/.bashrc
$ echo 'export PATH=$PATH:$ANDROID_HOME/emulator:$ANDROID_HOME/platform-tools' >> ~/.bashrc
$ source ~/.bashrc
# both commands should work without errors
$ java -version
$ adb devices
\end{lstlisting}

\noindent For macOS, the default SDK location is \q{\$HOME/Library/Android/sdk}, and
\q{JAVA\_HOME} can be set by \q{\$(/usr/libexec/java\_home -v 17)}.


\subsection{Installing Packages}
\markboth{Prototyping}{Installing Packages}

There are two ways to start with NativePHP. One option is to generate a new application from the
starter kit, which already contains NativePHP and its UI components.

\begin{lstlisting}[language=terminal]
$ laravel new fantasyland --using=nativephp/mobile-starter --no-node
\end{lstlisting}

\noindent Let's start our project as a regular Laravel application and add NativePHP on top of
it:

\begin{lstlisting}[language=terminal]
$ composer create-project laravel/laravel .
$ cp .env.example .env
$ php artisan key:generate
$ composer require nativephp/mobile nativephp/mobile-ui
# publish NativeServiceProvider (first time only)
$ php artisan vendor:publish --tag=nativephp-plugins-provider
# add the UI plugin into the allow-list
$ php artisan native:plugin:register nativephp/mobile-ui
$ php artisan native:plugin:list
\end{lstlisting}

\noindent Plugins that are missing in \q{NativeServiceProvider::plugins()} won't be included in
a build, which will result in blank screens for native screens of \q{nativephp/mobile-ui}.
Before running the installer, the application identity must be declared in the \q{.env}-file.
The \q{NATIVEPHP\_APP\_ID} is mandatory and follows the reverse domain notation. It becomes the
package name on Google Play and the bundle identifier on the App Store. Therefore, it is
difficult to change after the first release:

\begin{lstlisting}[language=terminal]
NATIVEPHP_APP_ID=com.tercad.fantasyland
NATIVEPHP_APP_VERSION=DEBUG
NATIVEPHP_APP_VERSION_CODE=1
\end{lstlisting}

\noindent The DEBUG version causes the device to refresh the Laravel bundle with each launch.
With a fixed version (such as 1.0.0), the \q{NATIVEPHP\_APP\_VERSION\_CODE} must be increased
before each build, otherwise the installed application may display old screens. Now, the native
part can be installed (\cref{img:fs-install}):

\begin{lstlisting}[language=terminal]
$ php artisan native:install
$ npm install
$ npm run build
\end{lstlisting}

\noindent The installer creates \q{config/nativephp.php}, the \q{native}-helper (so
\q{php native run} is the same as \q{php artisan native:run}), and the \q{nativephp}-directory
with Android and iOS projects and downloaded PHP binaries. It also asks whether to use PHP
binaries with the ICU library (\q{--with-icu}), which is needed only for the \q{intl}-extension
(Filament, for example) and adds about 30 MB to the Android build.

\img{img/native-install}{Installation of NativePHP for Android}{img:fs-install}

\noindent The \q{nativephp}-directory is ignored by Git and treated as ephemeral; it's
regenerated by \q{native:install}. Therefore, any manual fixes inside it will be lost after the
next reinstallation. That's why our repository stores such fixes in a separate
\q{\_patches}-folder.


\subsection{Starting Application}
\markboth{Prototyping}{Starting Application}

Before moving on to the device, it's worth checking that the Laravel part works at all. The same
application can be opened in a desktop browser:

\begin{lstlisting}[language=terminal]
$ php artisan serve
   INFO  Server running on [http://127.0.0.1:8000].
\end{lstlisting}

\noindent It's a quick way to iterate on routes, controllers, and Blade templates. However, the
browser doesn't have the native layer. The tab bar, toasts, and other features are only executed
inside the mobile shell.


\subsection{Running on Emulator}
\markboth{Prototyping}{Running on Emulator}

A virtual device is created in Android Studio (Device Manager $\rightarrow$ Create Virtual
Device). After that, NativePHP can list and start emulators, build the application, and install
it:

\begin{lstlisting}[language=terminal]
$ php artisan native:emulator android
$ php artisan native:run android
# rebuild, install, and push further changes
$ php artisan native:run android --watch
# stream Laravel logs from the Android application
$ php artisan native:tail
\end{lstlisting}

\noindent The \q{--watch}-option (as well as a separate \q{native:watch}-command) pushes changed
files from the folders listed in \q{config/nativephp.php} (\q{hot\_reload.watch\_paths}) and
reloads the current screen. For iOS, the same commands are executed with the \q{ios}-argument on
macOS, and the build is launched in the simulator.

\paragraph{Remarks} A first run rarely goes smoothly. Here are the problems that appeared on our
way, with their solutions:

\begin{itemize}
  \item \textbf{SDK location is not found}. Gradle failed with "\emph{sdk.dir property in
  local.properties file. Problem: Set with empty value}". The \q{ANDROID\_HOME}-variable was not
  exported in the current shell. Exporting it (see above) and, if needed, setting
  \q{sdk.dir=/home/<user>/Android/Sdk} in \q{nativephp/android/local.properties} solves it.

  \item \textbf{No matching ABIs}. The installation into the emulator failed with
  \q{INSTALL\_FAILED\_NO\_MATCHING\_ABIS}. NativePHP ships PHP binaries only for \q{arm64-v8a},
  while the default Android emulator on a PC is \q{x86\_64}. An ARM64 system image doesn't help
  either: the emulator stops with "\emph{Avd's CPU Architecture 'arm64' is not supported by the
  QEMU2 emulator on x86\_64 host}". So, on x86\_64 machines, the application has to be checked
  on a real device (with enabled USB debugging, \q{adb devices} lists it), while an Apple
  silicon Mac runs ARM64 emulators natively.

  \item \textbf{Laravel bundle is not extracted}. The application started with "\emph{Failed
  opening required '.../android/native.php'}". The Android shell didn't extract the archive into
  the storage. We had to patch \q{LaravelEnvironment.kt} with three fixes:
  (1) Changed from \q{context.getDir("storage", Context.MODE\_PRIVATE)} to \q{File(context.
  dataDir, "app\_storage")} to correctly use the app's data directory, (2) refactored the unzip
  function to stream directly from ZIP to disk instead of loading entire files into memory
  (which could fail with large bundles), and (3) modified the extraction
  to continue even when the version can't be read, as long as Laravel hasn't been extracted yet.
  Since the \q{nativephp}-directory is regenerated, such a fix is lost on reinstallation, so
  it's worth reporting it upstream.

  \item \textbf{Request is null}. Laravel failed during the boot with
  "\emph{UrlGenerator::\_\_construct(): Argument \#2 (\$request) must be of type
  Illuminate\textbackslash Http\textbackslash Request, null given}". The trace pointed to
  \q{URL::forceHttps()} in \q{AppServiceProvider::boot()}, which was added by the starter
  template.  On a device, requests don't come through a web server, and the URL generator can't
  be created at that stage. The first attempt was to call it only in production
  (\q{if (\$this->app->isProduction())}), but, in the end, the call was removed: the embedded
  runtime doesn't need HTTPS redirects at all.

  \item \textbf{Storage directory is missing on iOS}. Our cookie jar was created in
  \q{storage\_path('app/cookie.txt')}, which worked on Android but failed on iOS: the file
  couldn't be created there, since its folder wasn't guaranteed to exist. Don't rely on the
  folder structure of the Laravel skeleton on a device; create the directory before writing:

\begin{lstlisting}
$this->fcurl = storage_path('cookie.txt');
$dir = dirname($this->fcurl);
if (!is_dir($dir)) {
  mkdir($dir, 0777, true);
}
\end{lstlisting}

  \item \textbf{Mysterious NUL-file}. The very first commit contained a file named \q{NUL} with
  200 lines of "\emph{error: no devices/emulators found}". Most likely, it was a Windows habit to
  suppress output by \q{> NUL}, which on Linux creates a regular file. The real message was
  useful, though: \q{adb} didn't see any device. The file was deleted and added into
  \q{.gitignore}.
\end{itemize}

\noindent Finally, \q{php artisan native:debug} collects the environment details (PHP, NativePHP,
Java, Gradle, and Xcode on macOS) into one report. It's the first thing to attach to an issue
when something goes wrong.


\subsection{Keeping Framework Up to Date}
\markboth{Prototyping}{Keeping Framework Up to Date}

NativePHP for Mobile evolves quickly. During the first eight months of the project, it was
upgraded three times. Each upgrade was treated as a separate task since almost all of them
changed the project structure.

As a tip: the lock file is created by \q{native:install} and should be committed. It ensures
that every developer and the CI pipeline download the same PHP binaries:

\begin{lstlisting}[language=java]
{
  "php": {
    "version": "8.4.26",
    "icu": false
  }
}
\end{lstlisting}

\noindent The upgrade itself is short, while the follow-up is not:

\begin{lstlisting}[language=terminal]
$ composer update
# re-create native projects and download new PHP binaries
$ php artisan native:install --force
$ php artisan native:plugin:list
$ php artisan test
\end{lstlisting}

\noindent And, finally, the first screen is on the device... let's continue with the prototype!
